%%\documentclass[final,3p,times,twocolumn]{elsarticle}
\documentclass[final,3p,times,twocolumn]{report}

%% Use the option review to obtain double line spacing
%% \documentclass[preprint,review,12pt]{elsarticle}

%% Use the options 1p,twocolumn; 3p; 3p,twocolumn; 5p; or 5p,twocolumn
%% for a journal layout:
%% \documentclass[final,1p,times]{elsarticle}
%% \documentclass[final,1p,times,twocolumn]{elsarticle}
%% \documentclass[final,3p,times]{elsarticle}
%% \documentclass[final,3p,times,twocolumn]{elsarticle}
%% \documentclass[final,5p,times]{elsarticle}
%% \documentclass[final,5p,times,twocolumn]{elsarticle}

%% if you use PostScript figures in your article
%% use the graphics package for simple commands
%% \usepackage{graphics}
%% or use the graphicx package for more complicated commands
%% \usepackage{graphicx}
%% or use the epsfig package if you prefer to use the old commands
%% \usepackage{epsfig}

%% The amssymb package provides various useful mathematical symbols
\usepackage{amssymb}
%% The amsthm package provides extended theorem environments
%% \usepackage{amsthm}
%% The bm package lets you access bold symbols in math mode using the \boldsymbol command (useful to get bold greek letters).
\usepackage{bm}
%% The bbm package is contains the indicator function symbol \mathbbm{1}
\usepackage{bbm}
%% The amsmath package contains the split environment, letting you split equations into multiple lines.
%% See "https://www.sharelatex.com/learn/Aligning_equations_with_amsmath " for an explanation.
\usepackage{amsmath}
%% The lineno packages adds line numbers. Start line numbering with
%% \begin{linenumbers}, end it with \end{linenumbers}. Or switch it on
%% for the whole article with \linenumbers after \end{frontmatter}.
%% \usepackage{lineno}
%% The algorithm defines the algorithm floating environment and the algpseudocode package is useful for constructing Pseudo code.
\usepackage{algorithm}
\usepackage{algpseudocode}
%% For creating pictures
\usepackage{tikz}
%% For making algorithms float
\usepackage{float}
\newfloat{algorithm}{t}{lop}
%% For creating draft watermark
%\usepackage{draftwatermark}
%\SetWatermarkText{DRAFT}
%\SetWatermarkScale{1}

%% Declaring \argmin and \argmax operators:
\DeclareMathOperator*{\argmin}{arg\,min}
\DeclareMathOperator*{\argmax}{arg\,max}
%% Declare trace operator \Tr:
\DeclareMathOperator*{\Tr}{Tr}
%% Declare pdf functions
\DeclareMathOperator*{\Cat}{Cat}
\DeclareMathOperator*{\Dir}{Dir}
\DeclareMathOperator*{\DP}{DP}
\DeclareMathOperator*{\Beta}{Beta}
\DeclareMathOperator*{\GEM}{GEM}
\DeclareMathOperator*{\Stick}{Stick}
\DeclareMathOperator*{\Uniform}{Uniform}
%% shorthand for \boldsymbol and \overline
\let\bs\boldsymbol
\let\ol\overline
%% command that allows equations to be split across pages
%\allowdisplaybreaks[3]

%% natbib.sty is loaded by default. However, natbib options can be
%% provided with \biboptions{...} command. Following options are
%% valid:

%%   round  -  round parentheses are used (default)
%%   square -  square brackets are used   [option]
%%   curly  -  curly braces are used      {option}
%%   angle  -  angle brackets are used    <option>
%%   semicolon  -  multiple citations separated by semi-colon
%%   colon  - same as semicolon, an earlier confusion
%%   comma  -  separated by comma
%%   numbers-  selects numerical citations
%%   super  -  numerical citations as superscripts
%%   sort   -  sorts multiple citations according to order in ref. list
%%   sort&compress   -  like sort, but also compresses numerical citations
%%   compress - compresses without sorting
%%
%% \biboptions{comma,round}

% \biboptions{}


%\journal{MPhil in Scientific Computing}

\begin{document}

%% for elsarticle class
%\begin{frontmatter}

%% Title, authors and addresses

%% use the tnoteref command within \title for footnotes;
%% use the tnotetext command for the associated footnote;
%% use the fnref command within \author or \address for footnotes;
%% use the fntext command for the associated footnote;
%% use the corref command within \author for corresponding author footnotes;
%% use the cortext command for the associated footnote;
%% use the ead command for the email address,
%% and the form \ead[url] for the home page:
%%
%% \title{Title\tnoteref{label1}}
%% \tnotetext[label1]{}
%% \author{Name\corref{cor1}\fnref{label2}}
%% \ead{email address}
%% \ead[url]{home page}
%% \fntext[label2]{}
%% \cortext[cor1]{}
%% \address{Address\fnref{label3}}
%% \fntext[label3]{}

\begin{titlepage}
\begin{center}
\includegraphics[width=0.6\textwidth]{cam.png}\\

\vspace*{4cm}
\Huge\textbf{Compressive Sensing in Video Encoding}

\vspace*{1.5cm}
\LARGE
\textbf{Brian Azizi}
\vfill
Project Progress Report
%\normalsize
%A thesis submitted in partial fulfillment of the requirements for\\
%the degree of Master of Philosophy

\vspace{0.8cm}

\normalsize
Supervised by:\\
Dr Anita Faul

\vspace{0.2cm}
Department of Physics\\
Cavendish Laboratory\\
JJ Thomson Avenue\\
CB3 0HE\\
Cambridge, UK\\
\vspace{0.4cm}
\normalsize
\today

\end{center}
\end{titlepage}


\title{Compressive Sensing in Video Encoding}

%% use optional labels to link authors explicitly to addresses:
%% \author[label1,label2]{<author name>}
%% \address[label1]{<address>}
%% \address[label2]{<address>}

\author{Brian Azizi}

%\address{Cavendish Laboratory, Department of Physics, J J Thomson
 % Avenue, Cambridge. CB3 0HE}

\renewcommand{\abstractname}{Project Outcome}
\begin{abstract}                
The project outcome shall be a generic compressive sensing algorithm which takes as input 3D data which can be interpreted as video and encodes this efficiently and also reconstructs data when data is missing.
\end{abstract}

% % Acknowledgements
% \renewcommand{\abstractname}{Acknowledgements}
% \begin{abstract}
% Shoutouts to Nikos and Anita and Georgios and Foba.
% \end{abstract}
\tableofcontents

%\end{frontmatter}

%%
%% Start line numbering here if you want
%%
% \linenumbers

%% main text
\chapter{Problem Formulation}
\label{chap:Intro}
\section{Introduction}
Since beginning this project at the end of March, I have spent most of my time reviewing the work done by Georgios Pilikos and Anita Faul on Bayesian Compressive Sensing in image reconstruction.
The objective of my project is to extend their work to 3 dimensional signals (videos).

I will give a description of the compressive sensing (CS) framework and the theory behind the building blocks of their software (RVM, wavelet transforms, cascade).
I will then summarize my findings from reviewing the code base.
Finally, I will outline my current ideas on how to extend their work.

\section{Compressive Sensing}
Consider a digital signal $\bs x \in \mathbb{R}^N$. 
For example, $\bs x$ could be a 1-dimensional signal such as a speech signal.
In that case, the individual components of the vector $\bs x$ contain (typically equidistant) samples of the corresponding analogue (continuous) speech signal.

Conventionally, the sampling rate is based on the sampling theorem by Shannon and Nyquist.
The theorem states that it is possible to prefectly reconstruct the original analogue signal from the acquired samples if the sampling rate is at least twice the highest frequency in the signal.

In recent years, an emerging framework called compressive sensing (CS) has enabled us to perfectly reconstruct a certain class of signals from far fewer samples than required by the Shannon-Nyquist Theorem. 
The requirement is that the desired signal $\bs x$ exhibits \emph{sparsity} in some domain.
This means that $\bs x$ (or some basis transformation of it) has only a few non-zero elements.

CS allows for reconstruction of $\bs x \in \mathbb{R}^N$ using only $M$ samples, where $M << N$.
Let $\bs y \in \mathbb{R}^M$ be the vector containing these measurements.
Our aim is to find $\bs x$ such that
\begin{equation*}
\bs\Omega\bs x = \bs y
\end{equation*}
where $\bs\Omega$ is a known $M\times N$ matrix called the \emph{sensing matrix}.


\section{2D signals}
For illustrative purposes, we will regard $\bs x$ as an image.
In this case, $N$ is equal to the total number of pixels and we can treat $\bs x$ as a matrix whose dimensions are the resolution of the image that we have ``unrolled'' into a vector.
We would do this by putting all the rows of the matrix sequentially into a single row vector and transposing (so that we end up with a column vector).




\chapter{Sparse Bayesian Learning}
\label{chap:RVM}
\cite{tipping2001}
\chapter{Wavelets}
%\section*{Acknowledgements}
%Here I acknowledge the assistance of my supervisor, my industrial sponsor,
%and the effects of caffeine on my ability to produce this report on time.

%% The Appendices part is started with the command \appendix;
%% appendix sections are then done as normal sections
%\appendix

%% References
%%
%% Following citation commands can be used in the body text:
%% Usage of \cite is as follows:
%%   \cite{key}         ==>>  [#]
%%   \cite[chap. 2]{key} ==>> [#, chap. 2]
%%

%% References with bibTeX database:
%\section*{Bibliography}
\bibliographystyle{elsarticle-num}
%\bibliographystyle{apa}
\bibliography{referencesRVM.bib}

%% Authors are advised to submit their bibtex database files. They are
%% requested to list a bibtex style file in the manuscript if they do
%% not want to use elsarticle-num.bst.

%% References without bibTeX database:

% \begin{thebibliography}{00}

%% \bibitem must have the following form:
%%   \bibitem{key}...
%%

% \bibitem{}

% \end{thebibliography}


\end{document}

%%
%% End of file `rvm.tex'.